% Input by appendix.tex. Numbers: notes/E-evaluation.md, notes/appendix/P-pipeline.md.
\section{Time and Storage Analysis}
\label{app:analysis}

This appendix derives the time and storage cost of each operation from the structures
in Appendices~\ref{app:objects}--\ref{app:algorithms} and checks each bound against
measurements. Measurements come from the evaluation's campaigns
(Section~\ref{sec:eval}). Each is one sample unless noted; storage bytes are exact and
times are wall clock.

\paragraph{Notation.} A tree has $N$ inodes. A file has length $F$ and $C \approx
F/\bar{c}$ chunks, with measured mean chunk size $\bar{c} \approx 19$\,KB. Pages hold at
most $P = 8$\,KiB. The mapping tree of a file has height $h_F = \lceil \log_{64} C \rceil$
or less (3 for 500\,MiB), the inode table $h_N = O(\log_{50} N)$, and a directory with $D$
entries $h_D = O(\log D)$. A Commit captures a set $\Delta$ of changed inodes; changed
file $i$ received $w_i$ written bytes in $e_i$ separate edit runs. $T$ is the whole-file
cutoff (128\,KiB). $H$ is history length and $B$ the number of Branches.

\subsection{Time}

\begin{table}[h]
\centering
\small
\begin{tabular}{p{0.15\linewidth} p{0.30\linewidth} p{0.17\linewidth} p{0.28\linewidth}}
\toprule
Operation & Cost & Grows with & Measured (median) \\
\midrule
Mount & Template copy $+$ one root object & nothing & 17--43\,ms across machines \\
WRITE $w$ & $O(w + \log R)$ overlay rows; no Store I/O & $w$ & -- \\
READ $w$, cold & $O(w + h_F + h_N)$ pages; $\leq$ cap$+1$ records and $\leq 4T$ decoded per object & $w$; $\log F$, $\log N$ & 158\,MB/s large file; flat across history ($<$1\%) \\
Capture & Generation flip & nothing & -- \\
Commit, construction & $O\big(\sum_{i\in\Delta}(w_i + e_i h_F P) + |\Delta| h_N P\big)$ & changes only & 100\,B edit: 24.4 / 24.5 / 28.1\,ms at 10 / 100 / 500\,MiB \\
Commit, lane & 2 head reads, $\leq$1,024-inode merge, one transaction & nothing ($N$, $H$, $B$) & no-change Commit 15--21\,ms \\
Commit, new data & $O(w)$ chunking, hashing, zstd, packing & $w$ & 500\,MiB: 1.36\,s ($\approx$385\,MB/s) \\
Fork & One Branch row & nothing & -- \\
Merge & Pages whose child ids differ, windows of 256 serials; $\leq$1\,MiB byte compare per text file & changed serials & -- \\
Unmount & Drain, close, delete overlay file & open state & 14--31\,ms \\
\bottomrule
\end{tabular}
\caption{Time per operation. Mount, Commit and unmount times are CLI calls including
process launch and routing, which costs 26--54\,ms alone; the Commit times are cold.}
\label{tab:time}
\end{table}

\paragraph{File size does not enter a Commit.} A localized edit touches one
root-to-boundary path per edit boundary, so its construction cost is $O(w + e\,h_F P)$,
and $h_F$ grows by one level per 64--128$\times$ growth in $F$. A 50$\times$ increase in
file size (10 to 500\,MiB) raised the cold Commit of a 100\,B overwrite from 24.4 to
28.1\,ms; a copy-up overlay grew from 191 to 604\,ms over the same range, and its time is
$\Theta(F)$. In a unit test the same three-extent overwrite read 3 mapping pages on an
8\,MB file and 4 on a 24\,MB file.

\paragraph{Tree size, history and Branch count do not enter the lane.} The publication
lane holds two head reads, a merge bounded to 1,024 inodes, and one transaction that
inserts a Commit and advances one Branch head. None of these reads the tree beyond the
bounded merge, walks history or enumerates Branches. Work that does grow---capture of
large changes, construction, compression, pack writes, and a full merge when the head
has moved---runs before the lane, so concurrent Commits overlap in all of it
(Figure~\ref{fig:twocommits}).

\paragraph{Reads are bounded by structure, not history.} A cold read of $w$ bytes costs
$O(w)$ plus one root-to-leaf path in the inode table and in the file's mapping tree. Each
stored object resolves a delta chain of at most 16 (whole file) or 8 (chunk) records,
decoding at most $4T = 512$\,KiB and reading at most $2T = 256$\,KiB encoded; a chain
that would exceed these bounds is re-based onto its root when written. Reads of a 65\,MB
tree and of a 500\,MiB file were within 1\% of each other before and after 40 and 5
Commits of history. Large-file reads are slow in absolute terms (158\,MB/s cold against
943\,MB/s for ovnfs) because each 128\,KiB FUSE request is a round trip with one
background request outstanding; that is a kernel-route cost, not a history cost.

\paragraph{Per-call overhead.} For per-tool-call \ws{}s, the cost added to each call is
mount, Commit and unmount. On a Docker machine with a SQLite Store the clean tiny-workload
row measured 19.2, 16.0 and 16.8\,ms, about 52\,ms in all with process launch included.
Sandbox creation is separate and amortized over many calls: 0.39\,s median on Docker,
44--68\,ms natively on macOS and in Apple containers. File calls that need no process
pay neither, because they run in the SDK's process (Section~\ref{sec:filecalls}).

\paragraph{Where time is lost.} New data is written twice, once into the overlay and
again into the Store, so a 500\,MiB random write committed in 1.36\,s against 1.0\,s for
a copy-up overlay. FUSE publication of every \texttt{write(2)} makes write bandwidth a
small fraction of a kernel overlay's (Appendix~\ref{app:modal}). Large files saved under
a new inode are re-chunked in full at Commit.

\subsection{Storage}

\begin{figure}[h]
\centering
\resizebox{\linewidth}{!}{% Measured time and storage. Input by appendix-analysis.tex.
\begin{tikzpicture}
\begin{groupplot}[group style={group size=3 by 1, horizontal sep=1.55cm},
  width=5.9cm, height=4.6cm,
  tick label style={font=\tiny}, label style={font=\scriptsize}, title style={font=\scriptsize\bfseries},
  legend style={font=\tiny, draw=none, fill=none}, ymode=log, grid=major, grid style={gray!20}]
% (a) Commit time vs file size
\nextgroupplot[title={(a) Cold Commit, 100\,B overwrite}, xmode=log, xlabel={file size (MiB)}, ylabel={ms},
  xtick={10,100,500}, xticklabels={10,100,500}, ymin=10, ymax=1000, legend pos=north west]
\addplot[mark=square*, orange!80!black, thick] coordinates {(10,190.9) (100,221.1) (500,603.9)};
\addlegendentry{copy-up (ovnfs)}
\addplot[mark=*, blue!70!black, thick] coordinates {(10,24.4) (100,24.5) (500,28.1)};
\addlegendentry{\sys{}}
% (b) store growth per Commit
\nextgroupplot[title={(b) Store growth per Commit}, ybar, bar width=4pt, ylabel={bytes},
  symbolic x coords={L10,L100,L500,I10,A5,DEL500}, xtick=data, x tick label style={rotate=40, anchor=east, font=\tiny},
  ymin=1e3, ymax=1e12, legend pos=north west, enlarge x limits=0.12]
\addplot[fill=orange!60, draw=orange!80!black] coordinates {(L10,10.5e6) (L100,104.9e6) (L500,524.3e6) (I10,10.5e6) (A5,613e3) (DEL500,524.3e6)};
\addlegendentry{copy-up (ovnfs)}
\addplot[fill=blue!40, draw=blue!70!black] coordinates {(L10,5.5e3) (L100,5.3e3) (L500,7.4e3) (I10,5.3e3) (A5,18.4e3) (DEL500,8.2e3)};
\addlegendentry{\sys{}}
% (c) logical / stored
\nextgroupplot[title={(c) Logical $\div$ stored bytes}, xbar, bar width=5pt, xmode=log, ymode=normal,
  symbolic y coords={M01 tiny files,M08 SQLite DB,D10 repo replay,M02 tree copies,D03 repo replay,M03 text edits,H05 20 projects,H01 500 Commits},
  ytick=data, y tick label style={font=\tiny}, xmin=0.5, xmax=4000, xlabel={ratio (SQLite Store)},
  nodes near coords, nodes near coords style={font=\tiny, anchor=west}, point meta=explicit symbolic,
  enlarge y limits=0.08]
\addplot[fill=blue!40, draw=blue!70!black] coordinates {
 (0.91,M01 tiny files) [0.91] (13.73,M08 SQLite DB) [13.7] (11.71,D10 repo replay) [11.7] (16.97,M02 tree copies) [17.0]
 (27.48,D03 repo replay) [27.5] (29.27,M03 text edits) [29.3] (275.5,H05 20 projects) [276] (1225.9,H01 500 Commits) [1226]};
\end{groupplot}
\end{tikzpicture}
}
\caption{Measured time and storage. (a) Cold Commit time of a 100\,B overwrite is flat in file size for \sys{} and grows with it for a copy-up overlay. (b) Store growth per Commit for large-file edits (log scale; DEL500 truncates 500\,MiB). (c) Logical snapshot bytes over stored bytes across workloads; tiny files are the one case where the Store is larger than the data.}
\label{fig:charts}
\end{figure}

\paragraph{Cost of one change.} Store growth from a Commit is new payload plus new
metadata plus pack and locator overhead. Unchanged files, chunks, pages and subtrees
cost nothing, because identities are content hashes and updates rewrite only paths to
changes.
\begin{itemize}[leftmargin=*,nosep]
\item \emph{Payload.} A localized edit of $w$ bytes in a chunked file creates between
$\lceil w/32\,\mathrm{KiB}\rceil$ and $\lceil w/8\,\mathrm{KiB}\rceil + 1$ new chunks,
because only the replacement is chunked and boundary chunks become slices of existing
chunks. Payload is therefore at most $w$ plus per-record framing (or a smaller delta),
independent of $F$. A 100\,B random overwrite stored 153\,B of payload at 10, 100 and
500\,MiB. A small file ($< T$) is rewritten whole, so its payload is a delta against
the previous version, bounded by its compressed size.
\item \emph{Metadata.} A changed file rewrites at most $h_F$ mapping pages, a 93\,B file
state, $h_N$ inode-table pages and a 116\,B root; a name change also rewrites $h_D$
directory pages. Canonically this is $O((h_F + h_N + h_D)P)$, 17--26\,KB for a
mid-file overwrite of a large file. Each rewritten page is stored as a COPY/INSERT
program against the page it replaces, so 3.7--5.7\,KB is stored. Across mixed histories
a least-squares fit gives 414\,B per edit, 526\,B per append, 464\,B per rename,
1,111\,B per delete and 1,360\,B per new file.
\item \emph{Store overhead.} SQLite pages, locators and indexes add 1--20\% over the
stored records (1.01--1.20\X{}; PostgreSQL 1.05--1.39\X{}). Per object the floor is
about 200\,B, which dominates only for tiny files.
\end{itemize}

\begin{table}[h]
\centering
\small
\begin{tabular}{p{0.32\linewidth} p{0.20\linewidth} p{0.38\linewidth}}
\toprule
Change & Cost bound & Measured \\
\midrule
100\,B overwrite of a 500\,MiB file & $O(w + h_F P)$ & 7.4\,KB growth per Commit (copy-up overlay: 524.3\,MB) \\
16\,KiB append to a 5\,MiB file & $O(w)$ & 18.4\,KB per Commit (613\,KB) \\
Truncate 500\,MiB off a 1,000\,MiB file & $O(h_F P)$ & 8.2\,KB (524.3\,MB) \\
Remove a 500\,MiB file & $O(h_N P)$ & 24.6\,KB; frees nothing \\
Copy a 50\,MB tree 10 times & $O(\text{metadata})$ & 15.9\,KB per copy Commit, 0 payload \\
Fork a Branch & one row; first Commit pays its paths & 5.4\,KB per identical Branch \\
Second project in a shared Store & $O(\text{metadata})$ for shared content & 220\,KB per Init (first: 23.6\,MB) \\
Commit in a 65\,MB tree (5 mixed ops) & $O(\sum$ changes$)$ & 5.9\,KB growth per Commit over 500 Commits \\
\bottomrule
\end{tabular}
\caption{Store growth per change. Copy-up figures are ovnfs on the same host.}
\label{tab:storagechange}
\end{table}

\paragraph{Cost of a history.} Store size after $k$ Commits is the initial tree plus
the sum of per-change costs; it does not grow with the size of each snapshot. Over 500
Commits on a 65\,MB tree (32.5\,GB of logical snapshots) the Store reached 26.5\,MB:
23.6\,MB of initial tree plus 5.9\,KB per Commit, a ratio of 1,226\X{} between logical
and stored bytes. Table~\ref{tab:ksnap} compares the growth of $k$ snapshots in which
each snapshot changes $m$ files of size $F$ by $w$ bytes. Copy-up overlays and
whole-file snapshot images grow with $F$; block copy-on-write grows with the touched
blocks, but has no cross-file deduplication or deltas; \sys{} grows with $w$ plus a
logarithmic path. A fully repacked Git pack is the tightest point of reference---1.10 to
1.31\X{} smaller on the measured histories, because one global repack deltas across all
versions---but Git keeps no live mountable view and must repack to reach that size.

\begin{table}[h]
\centering
\small
\begin{tabular}{p{0.30\linewidth} p{0.27\linewidth} p{0.33\linewidth}}
\toprule
System & Growth for $k$ snapshots & Measured point \\
\midrule
Copy-up overlay; whole-file snapshot images & $k\,m\,F$ & 100\,B edits to 100\,MiB file: 1.15\,GB after 10 \\
Block copy-on-write (Btrfs, ImageFS-blk) & $k\,m\,\lceil w/b\rceil b$ $+$ metadata & 4\,KiB change to 512\,MiB file: about 36\,KiB (Btrfs; recovered report) \\
Git after full repack & $\approx k\,m\,\delta(w)$, after repack & 1.10--1.31\X{} smaller than \sys{} \\
\sys{} & $k\,(m(w + c\,h_F P') + h_N P')$ & same 10 edits: 106.3\,MB, of which 104.9\,MB is the initial file \\
\bottomrule
\end{tabular}
\caption{Growth over $k$ snapshots. $b$ is the block size, $\delta(w)$ a delta of $w$
bytes, and $P'$ the stored size of a page program.}
\label{tab:ksnap}
\end{table}

\paragraph{Trade-offs taken.} Several storage choices were made, or refused, on their
time cost, following the rule that a large speed loss is not traded for a small storage
gain (Table~\ref{tab:tradeoffs}).

\begin{table}[h]
\centering
\small
\begin{tabular}{p{0.30\linewidth} p{0.30\linewidth} p{0.30\linewidth}}
\toprule
Choice & Time effect & Storage effect \\
\midrule
Re-base at the chain bound & $+$22.5\% mean decode; max 520,575\,B $\leq$ 4$T$ budget & Edit Commit payload 290\,KB $\rightarrow$ 120\,KB \\
Losing-base back-off & 500\,MiB append: 26,966 $\rightarrow$ 22 trials, 7.9 $\rightarrow$ 5.5\,s & Identical bytes \\
Win saving $<$1/8 counts as loss & 100\,MiB log append 1,540 $\rightarrow$ 977\,ms & $+$1.04\% stored \\
zstd-3, not zstd-6 (refused) & zstd-6 costs 2.6\X{} compression CPU & zstd-6 would save 4.4\% \\
Per-object records, no group compression (refused) & Group: decode 256\,KiB to serve 2\,KB & Group would save 22\% \\
Per-chunk frames & Random-access reads & Text 1.25\X{} larger than whole-file zstd \\
No garbage collection or repack & No background work & Deleted data is retained \\
\bottomrule
\end{tabular}
\caption{Speed and storage trade-offs, from the design records.}
\label{tab:tradeoffs}
\end{table}
